% Image credits for the photographs and AI illustrations of Book 5 (University, Year 3), Hindi edition.
% Machine-readable license records live in images/book5/CREDITS.md; keep this
% file in sync with frontmatter/image-credits-book5.tex (one \item per item there).
% Names of people, institutions, collections, works, place names and licences are
% kept exactly as in the English credits (Latin script); only the frame and the
% descriptions are translated.
\chapter*{चित्रों का श्रेय}

इस पुस्तक के रेखाचित्र, मानचित्र और चार्ट मौलिक हैं; प्रकाशित आँकड़ों से बनाए गए मानचित्र और
चार्ट अपना स्रोत चित्र-परिचय में या इस सूची में बताते हैं। फ़ोटो उनके अपने-अपने मुक्त लाइसेंसों के
अंतर्गत, उनके लेखकों के प्रति आभार के साथ, इस्तेमाल किए गए हैं; हर फ़ोटो के पूरे स्रोत-लिंक और
लाइसेंस के पते पुस्तक के भंडार में \texttt{images/book5/CREDITS.md} में दिए गए हैं।

\bigskip
फ़ोटो और मौलिक रेखाचित्रों के साथ-साथ, कुछ चित्र AI से बनाए गए चित्रण हैं, जो पुस्तक के
संपादकों के लिखे निर्देशों से OpenAI के चित्र-निर्माण से बनाए गए, और शामिल करने से पहले उनकी
वैज्ञानिक शुद्धता जाँची गई। वे केवल सामान्य दृश्य दिखाते हैं, कभी कोई असली जगह, असली व्यक्ति
या असली नमूना नहीं। हर बनाए गए चित्र का पूरा निर्देश भंडार में
\texttt{images/book5/ai/PROMPTS.md} में दर्ज है।

\bigskip
\noindent फ़ोटो और आँकड़े, अध्याय के अनुसार:
\begin{itemize}\small\setlength{\emergencystretch}{2em}
  \item अध्याय~1 (भूकंपविज्ञान और पृथ्वी का प्रतिबिंबन): PREM, A.~M.~Dziewonski और D.~L.~Anderson (1981), ObsPy में सारणीबद्ध रूप में; ak135, B.~L.~N.~Kennett, E.~R.~Engdahl और R.~Buland (1995)।
  \item अध्याय~3 (गुरुत्व और पृथ्वी की आकृति): जुड़वाँ GRACE Follow-On उपग्रह, कलाकार की परिकल्पना, NASA/JPL-Caltech, सार्वजनिक डोमेन (Wikimedia Commons)। हिमचादरों का द्रव्यमान संतुलन IMBIE Team (2021), UK Polar Data Centre, Open Government Licence v3.0 से।
  \item अध्याय~4 (भू-चुंबकत्व और पुराचुंबकत्व): International Geomagnetic Reference Field, चौदहवीं पीढ़ी, IAGA Division~V Working Group V-MOD, गुणांक NOAA NCEI द्वारा वितरित।
  \item अध्याय~5 (स्थलमंडल के ऊष्मीय मॉडल): महासागर तल की आयु, M.~Seton और अन्य (2020), EarthByte, CC BY 4.0; ETOPO1 वैश्विक उच्चावच, NOAA NCEI, सार्वजनिक डोमेन।
  \item अध्याय~6 (प्रावार संवहन और प्लेट-चालक बल): Ballyhillin, Malin Head के उठे हुए पुलिन, A.~F.~Borchert, CC BY-SA 4.0 (Wikimedia Commons); PREM, A.~M.~Dziewonski और D.~L.~Anderson (1981)।
  \item अध्याय~7 (विरूपण की क्रियाविधियाँ और अपरूपण क्षेत्र): ऑगन मायलोनाइट, Mikenorton, CC BY-SA 4.0 (Wikimedia Commons)।
  \item अध्याय~9 (विस्तारी और स्ट्राइक-स्लिप विवर्तनिकी): Death Valley के दक्षिण में बेसिन-एंड-रेंज स्थलाकृति, Ken Lund, CC BY-SA 2.0 (Wikimedia Commons)।
  \item अध्याय~14 (स्थिर-समस्थानिक भू-रसायन): LR04 नितलस्थ $\delta^{18}$O स्टैक, L.~E.~Lisiecki और M.~E.~Raymo (2005)।
  \item अध्याय~15 (जल--चट्टान अंतर्क्रिया): Palette Spring, Mammoth Hot Springs की ट्रैवर्टीन वेदिकाएँ, Jim Peaco, National Park Service, सार्वजनिक डोमेन (Wikimedia Commons)। साम्य स्थिरांक PHREEQC डेटाबेस phreeqc.dat से, D.~L.~Parkhurst और C.~A.~J.~Appelo, U.S. Geological Survey, सार्वजनिक डोमेन।
  \item अध्याय~16 (अनुक्रम स्तरिकी): Helper, Utah के ऊपर Book Cliffs, David Mackie, सार्वजनिक डोमेन (Wikimedia Commons); LR04 नितलस्थ $\delta^{18}$O स्टैक, L.~E.~Lisiecki और M.~E.~Raymo (2005)।
  \item अध्याय~17 (संलक्षणी मॉडल): Arisaig Formation में टीलेदार तिर्यक स्तरण, Michael C.~Rygel, CC BY-SA 3.0 (Wikimedia Commons), काटा हुआ।
  \item अध्याय~18 (मात्रात्मक भू-आकृतिविज्ञान): Niagara Gorge और Horseshoe Falls, Ken Lund, CC BY-SA 2.0 (Wikimedia Commons), काटा हुआ।
  \item अध्याय~19 (भूभौतिकीय तरल गतिकी): NCEP--NCAR Reanalysis~1 की आंचलिक-माध्य पवनें, आँकड़े NOAA PSL, Boulder, Colorado, USA द्वारा प्रदत्त; World Ocean Atlas 2023, NOAA NCEI।
  \item अध्याय~20 (वायुमंडलीय गतिकी): Iceland के निकट एक निम्न-दाब तंत्र, NASA/GSFC, MODIS Rapid Response Team, Jacques Descloitres, सार्वजनिक डोमेन (Wikimedia Commons); NCEP--NCAR Reanalysis~1 की आंचलिक-माध्य पवनें, आँकड़े NOAA PSL, Boulder, Colorado, USA द्वारा प्रदत्त।
  \item अध्याय~21 (महासागरीय गतिकी): अनुसंधान पोत से छोड़ा जा रहा एक Argo प्लव, NOAA Global Ocean Monitoring and Observing Program, सार्वजनिक डोमेन (Wikimedia Commons)।
  \item अध्याय~22 (जलवायु मॉडल): NASA Center for Climate Simulation, Goddard Space Flight Center के सुपरकंप्यूटर, 2010, सार्वजनिक डोमेन (Wikimedia Commons); NCEP--NCAR Reanalysis~1 का सतही वायु तापमान और वायुमंडल के शीर्ष का विकिरण, आँकड़े NOAA PSL, Boulder, Colorado, USA द्वारा प्रदत्त।
  \item अध्याय~24 (हिमनदविज्ञान): दक्षिण-पश्चिमी Greenland का एक निर्गम हिमनद, NASA / Jim Yungel, Operation IceBridge, 2014, सार्वजनिक डोमेन (Wikimedia Commons); Greenland की सतही ऊँचाइयाँ ETOPO1 Ice Surface, NOAA NCEI से; आर्कटिक समुद्री-हिम विस्तार NSIDC Sea Ice Index, version~4 से (Fetterer और अन्य, 2025, doi:10.7265/a98x-0f50)।
  \item अध्याय~25 (सुदूर संवेदन और GIS): Three Sisters ज्वालामुखीय केंद्र, Oregon का 1995--2001 का व्यतिकरण-लेख, U.S. Geological Survey (C.~Wicks और अन्य, 2002), सार्वजनिक डोमेन (Wikimedia Commons); प्रयोगशाला स्पेक्ट्रम ECOSTRESS Spectral Library, version~1.0 से (S.~K.~Meerdink, S.~J.~Hook, D.~A.~Roberts और E.~A.~Abbott, 2019; A.~M.~Baldridge और अन्य, 2009), Jet Propulsion Laboratory, California Institute of Technology, Pasadena, California की सौजन्य से पुनर्प्रस्तुत।
  \item अध्याय~26 (ग्रहों के आंतरिक भाग और सतहें): Barringer Meteor Crater, Arizona, USGS/D.~Roddy, सार्वजनिक डोमेन (Wikimedia Commons); Olympus Mons, Viking ऑर्बिटर मोज़ेक, NASA/\allowbreak JPL/\allowbreak USGS, सार्वजनिक डोमेन (Wikimedia Commons)। क्रेटर उत्पादन और कालक्रम फलन G.~G.~Michael के Craterstats सॉफ़्टवेयर में दिए रूप में, G.~Neukum, B.~A.~Ivanov और W.~K.~Hartmann (2001) के अनुसार।
  \item अध्याय~27 (ग्रहीय वायुमंडल और वासयोग्यता): मंगल पर सूखी घाटियों का एक जाल, Mars Express HRSC, ESA/\allowbreak DLR/\allowbreak FU Berlin, CC BY-SA 3.0 IGO (Wikimedia Commons); Ligeia Mare, Titan, Cassini रडार मोज़ेक, NASA/\allowbreak JPL/\allowbreak USGS, सार्वजनिक डोमेन (Wikimedia Commons)। वासयोग्य-क्षेत्र फ़िट R.~K.~Kopparapu और अन्य (2013); मीथेन का वाष्प दाब NIST Chemistry WebBook से।
  \item अध्याय~28 (बहिर्ग्रह): James Webb Space Telescope द्वारा चित्रित HIP~65426~b, NASA/\allowbreak ESA/\allowbreak CSA, A~Carter (UCSC), ERS~1386 टीम, और A.~Pagan (STScI), CC BY~4.0 (Wikimedia Commons)। पुष्ट ग्रह NASA Exoplanet Archive की Planetary Systems Composite Parameters सारणी से, 23 अक्टूबर 2026 को देखी गई: इस शोध में NASA Exoplanet Archive का उपयोग हुआ है, जिसका संचालन California Institute of Technology करता है, National Aeronautics and Space Administration के साथ अनुबंध में, Exoplanet Exploration Program के अंतर्गत। द्रव्यमान--त्रिज्या के मॉडल वक्र L.~Zeng और अन्य (2016, 2019) के।
  \item अध्याय~29 (आकाशगंगा और मंदाकिनियाँ): GigaGalaxy Zoom परियोजना का आकाशगंगा पैनोरमा, ESO/S.~Brunier, CC BY~4.0 (Wikimedia Commons); सर्पिल मंदाकिनी M101, European Space Agency और NASA (Hubble, heic0602; Hubble आँकड़े K.~D.~Kuntz, F.~Bresolin, J.~Trauger, J.~Mould और Y.-H.~Chu; चित्र-संसाधन Davide De~Martin), CC BY~4.0 (Wikimedia Commons)। आकाशगंगा के वृत्तीय वेग A.-C.~Eilers, D.~W.~Hogg, H.-W.~Rix और M.~K.~Ness (2019) से; E.~Pouliasis, P.~Di~Matteo और M.~Haywood (2017) का द्रव्यमान मॉडल; Large Magellanic Cloud के OGLE क्लासिकी सीफ़ीड, I.~Soszy\'nski और अन्य (2015), Acta Astronomica 65, 297, CDS, Strasbourg द्वारा वितरित (VizieR J/AcA/65/297)।
  \item अध्याय~30 (तत्त्व कहाँ से आते हैं): Crab Nebula, NASA, ESA, J.~Hester और A.~Loll (Arizona State University), सार्वजनिक डोमेन (Wikimedia Commons); NGC~4993 में GW170817 का किलोनोवा, NASA और ESA, आभार A.~J.~Levan, N.~R.~Tanvir, A.~Fruchter और O.~Fox (Hubble, heic1717a), CC BY~4.0 (Wikimedia Commons)। सौर प्रचुरताएँ M.~Asplund, A.~M.~Amarsi और N.~Grevesse (2021) से; न्यूट्रॉन-प्रग्रहण अनुप्रस्थ काट KADoNiS v0.3 से (I.~Dillmann, R.~Plag, F.~K\"appeler और T.~Rauscher); अर्ध-आयु IAEA Live Chart of Nuclides (ENSDF) से; Geneva--Copenhagen सर्वेक्षण की धात्विकताएँ, L.~Casagrande और अन्य (2011) के पुनर्विश्लेषण में, CDS, Strasbourg द्वारा वितरित (VizieR J/A+A/530/A138)।
\end{itemize}
